\documentclass[12pt]{article}
\usepackage{amsmath, amssymb}

\title{CSE 400: Fundamentals of Probability in Computing\\
Lecture 13 (Tutorial) – Consolidated Exam Revision Notes}
\date{}
\begin{document}

\maketitle

These notes systematically organize all definitions, distributions, assumptions, identities, and standard probability results appearing in Tutorial-1 (Lecture 13). The presentation follows the mathematical structure underlying the problems and consolidates the tools required for solving them in an examination setting.

\section*{1. Combinatorial Structures and Multinomial Counting}

\subsection*{Partitioning Distinct Objects}

Given 20 distinct objects divided into four groups of five objects each, where:

Each group contains exactly 5 objects,

Order within each group is irrelevant,

Order among the groups is irrelevant,

the number of possible divisions is

\[
\frac{20!}{(5!)^4 \, 4!}.
\]

The factor $(5!)^4$ removes ordering within each group.

The factor $4!$ removes ordering among the four groups.

\subsection*{Probability via Counting}

If structural constraints are imposed (e.g., grouping by cuisine when exactly five objects of each type exist), probability is computed as

\[
\frac{\text{Number of favorable partitions}}{\text{Total number of partitions}}.
\]

\section*{2. Uniform Distribution on an Interval}

A passenger arrives at a time uniformly distributed between 7:00 and 7:30 A.M.

Let 
\[
X \sim \text{Uniform}(0,30).
\]

For $a < b$,

\[
P(a \le X \le b) = \frac{b - a}{30}.
\]

When arrivals (e.g., buses every 15 minutes) determine waiting time, probabilities are evaluated by measuring the corresponding subinterval lengths relative to the total interval length 30.

\section*{3. Conditional Probability and Basic Identities}

Let events $A$ and $B$.

\subsection*{Conditional Probability}

\[
P(A \mid B) = \frac{P(A \cap B)}{P(B)}, \quad P(B) > 0.
\]

\subsection*{Complement Rule}

\[
P(A^c) = 1 - P(A).
\]

\subsection*{Decomposition Identities}

\[
P(A \cap B^c) = P(A) - P(A \cap B).
\]

These relations are applied when probabilities of joint and marginal events are given.

\section*{4. Poisson Distribution}

\subsection*{Definition}

A random variable $X$ is Poisson with parameter $\lambda > 0$ if

\[
P(X = k) = \frac{e^{-\lambda} \lambda^k}{k!}, \quad k = 0,1,2,\dots
\]

\subsection*{Key Properties}

Mean: 
\[
E[X] = \lambda.
\]

Used for modeling counts over fixed intervals.

\subsection*{Cumulative and Complement Forms}

\[
P(X \ge m) = 1 - \sum_{k=0}^{m-1} \frac{e^{-\lambda} \lambda^k}{k!}
\]

\[
P(X \le m) = \sum_{k=0}^{m} \frac{e^{-\lambda} \lambda^k}{k!}
\]

\subsection*{Conditional Probability for Poisson}

\[
P(X \ge a \mid X \ge b) = \frac{P(X \ge a)}{P(X \ge b)}, \quad a \ge b.
\]

\section*{5. Exponential Series Identity}

The identity

\[
\sum_{i=0}^{\infty} \frac{\lambda^i}{i!} = e^{\lambda}
\]

is used for normalization of probability mass functions of the form

\[
p(i) = c \frac{\lambda^i}{i!}, \quad i = 0,1,2,\dots
\]

Imposing

\[
\sum_{i=0}^{\infty} p(i) = 1
\]

gives

\[
c = e^{-\lambda}.
\]

\section*{6. Binomial Distribution and Majority Rules}

\subsection*{Definition}

If each of $n$ independent trials succeeds with probability $p$, then

\[
X \sim \text{Binomial}(n,p)
\]

with

\[
P(X = k) = \binom{n}{k} p^k (1-p)^{n-k}.
\]

\subsection*{Majority Threshold}

For systems requiring at least half the components to function:

If $n = 5$, system operates if $X \ge 3$.

If $n = 3$, system operates if $X \ge 2$.

In general, if $n = 2k + 1$, operation requires $X \ge k + 1$.

Probability of operation:

\[
P(X \ge \text{threshold}) = \sum_{k=\text{threshold}}^{n} \binom{n}{k} p^k (1-p)^{n-k}.
\]

System comparisons are made by comparing these cumulative binomial probabilities.

\section*{7. Jury Decision Model}

Assume:

12 independent jurors,

Each makes the correct decision with probability $\theta$.

Let

\[
X \sim \text{Binomial}(12,\theta).
\]

A conviction requires at least 8 guilty votes:

\[
P(\text{correct decision}) = \sum_{k=8}^{12} \binom{12}{k} \theta^k (1-\theta)^{12-k}.
\]

Independence and identical probability assumptions are essential.

\section*{8. Normal Distribution and Q-Function}

Given density

\[
f_X(x) = \frac{1}{\sqrt{8\pi}} \exp\left(-\frac{(x+3)^2}{8}\right),
\]

this corresponds to a normal distribution with:

\[
\mu = -3, \quad \sigma^2 = 4.
\]

\subsection*{Standardization}

\[
Z = \frac{X - \mu}{\sigma}.
\]

\subsection*{Q-Function}

\[
Q(x) = P(Z > x).
\]

All required probabilities are expressed in terms of the Q-function after standardization.

Absolute value expressions are converted into interval probabilities and then standardized accordingly.

\section*{Core Mathematical Tools for Examination}

The tutorial relies on the following foundational structures:

Multinomial counting for equal partitions

Uniform distribution on finite intervals

Conditional probability and complement identities

Poisson distribution and cumulative complements

Exponential series normalization

Binomial distribution with majority thresholds

Normal distribution standardization and Q-function

These tools form the complete conceptual and computational framework required for solving all problems in Lecture 13 (Tutorial).

\end{document}